% SPDX-License-Identifier: GPL-3.0-or-later
% Copyright (C) 2026 Alessandro Greco
%% EPUB 无障碍替代文本。
\altfig{TCP 和 TLS 上的 Tor 封装示意图。外部观察者能看到 IP 地址、端口、流量大小、时间特征和 TLS 记录长度；没有 TLS 会话密钥则无法解析内部的 Tor cell。一个 TLS 记录可含多条 cell，也可能将一条 cell 分割到多个记录中，因此两者边界未必对齐。链路协议 v4 或更高版本中，固定长度 cell 为 514 字节，其中 5 字节为首部、509 字节为主体；还有可变长度 cell。只有 RELAY cell 的主体采用电路密钥逐层加密，PADDING 和 VPADDING 等 cell 不是 RELAY 消息。获得 TLS 密钥只能解析 cell，不能自动解密其 RELAY 主体。}
% 封装示意图：TLS 记录边界与 Tor cell 边界不必一致。
\begin{tikzpicture}[font=\sffamily\scriptsize,
  liv/.style={draw,rounded corners=2pt,minimum height=0.85cm,align=center},
  cel/.style={draw,fill=torviola!12,minimum height=0.7cm,minimum width=1.5cm,font=\sffamily\tiny}]
  \node[liv,fill=torgrigio!10,minimum width=10.6cm] (tcp) at (0,2.6) {TCP 字节流：分段长度可变};
  \node[liv,fill=torverde!12,minimum width=9.4cm] (tls) at (0,1.35) {TLS：记录与 cell 边界不对齐};
  \foreach \i/\x in {1/-3.6,2/-1.8,3/0,4/1.8,5/3.6} \node[cel] (c\i) at (\x,0) {514\,B};
  \node[anchor=west,font=\sffamily\tiny,text=torgrigio] at (4.6,0) {定长 v4+};
  \node[liv,fill=avviso!10,minimum width=4.6cm] (rel) at (-1.7,-1.4) {RELAY：509\,B 主体\\逐层加密};
  \node[liv,fill=torgrigio!8,minimum width=4.6cm] (pad) at (3.3,-1.4) {PADDING / VPADDING\\不是 RELAY};
  \draw[torgrigio] (c1.south west) -- (rel.north west); \draw[torgrigio] (c3.south east) -- (rel.north east);
  \node[anchor=west,align=left,font=\sffamily\tiny,text=torgrigio] at (5.8,2.6) {IP、端口、流量\\与时间可观察};
  \node[anchor=west,align=left,font=\sffamily\tiny,text=torgrigio] at (5.8,1.35) {TLS 记录长度\\在链路上可见};
  \node[anchor=west,align=left,font=\sffamily\tiny,text=torgrigio] at (5.8,0) {解密 TLS 后\\才能解析 cell};
  \node[anchor=west,align=left,font=\sffamily\tiny,text=torgrigio] at (5.8,-1.4) {RELAY 主体还需\\电路密钥};
  \node[font=\sffamily\tiny,text=torgrigio,align=center] at (0,-2.55) {仅为示意：TLS 记录与 cell 的边界不一定对齐。};
\end{tikzpicture}
